\documentclass[aspectratio=169, 10pt]{beamer}
% Shared Beamer setup for the three workshop decks.
% Build with: python tasks.py slides w1   (or compile slides/wN.tex from the slides folder)
%
% Build variants (set by tasks.py through latexmk -usepretex):
%   \finalbuild  hides PREP markers; tasks.py refuses a final build while any remain
%   \withnotes   shows speaker notes on a second screen
%   handout      beamer class option for the participant copy

\usetheme[progressbar=frametitle, numbering=fraction, block=fill]{metropolis}
\usepackage[T1]{fontenc}
\usepackage{booktabs, array, tabularx, pifont, xcolor, tikz, pgfplots}
\pgfplotsset{compat=1.18}
\usetikzlibrary{arrows.meta, positioning, fit, backgrounds, calc}
\usepackage[backend=biber, style=numeric-comp, sorting=none, maxbibnames=3, giveninits=true]{biblatex}
\addbibresource{../../../bibliography/references.bib}
\renewcommand*{\bibfont}{\tiny}
\setbeamertemplate{bibliography item}{\insertbiblabel}

\definecolor{okgreen}{HTML}{2E7D32}
\definecolor{badred}{HTML}{C62828}
\definecolor{warnamber}{HTML}{B26A00}
\definecolor{tierAcol}{HTML}{1B5E20}
\definecolor{tierBcol}{HTML}{0D47A1}
\definecolor{tierCcol}{HTML}{6A1B9A}
\definecolor{orgrecycler}{HTML}{00796B}
\definecolor{orglab}{HTML}{1565C0}
\definecolor{orgmaker}{HTML}{E65100}
\definecolor{orgbroker}{HTML}{C62828}
\definecolor{orgregulator}{HTML}{5D4037}
\setbeamercolor{alerted text}{fg=badred}

% Table marks used by generated results
\newcommand{\cmark}{\textcolor{okgreen}{\ding{51}}}
\newcommand{\xmark}{\textcolor{badred}{\ding{55}}}
\newcommand{\pmark}{\textcolor{warnamber}{\textbf{$\sim$}}}

% Evidence tiers (docs/evidence-policy.md): \ev{A}{key} prints a badge and a citation
\newcommand{\tier}[1]{\tikz[baseline=(t.base)]\node[fill=tier#1col, text=white, font=\tiny\bfseries,
  inner sep=1.1pt, rounded corners=1pt] (t) {#1};}
\newcommand{\ev}[2]{\tier{#1}\,\cite{#2}}
\newcommand{\pre}{{\tiny\itshape(preprint)}}
\newcommand{\grey}[1]{\tier{C}\,\cite{#1}\,{\tiny\itshape motivation and vocabulary only}}

% Preparation markers: visible while drafting, must be resolved before the final build
\newcommand{\prep}[1]{\ifdefined\finalbuild\else{\setlength{\fboxsep}{1.5pt}%
  \colorbox{warnamber!20}{\parbox{0.92\linewidth}{\tiny\textbf{PREP:} #1}}}\fi}

\ifdefined\withnotes
  \setbeameroption{show notes on second screen=right}
\fi

% Generated experiment output, with a visible fallback when it has not been run
\newcommand{\result}[1]{\IfFileExists{../results/#1.tex}{\input{../results/#1.tex}}%
  {\colorbox{badred!15}{\scriptsize Missing result \texttt{\detokenize{#1}}: run \texttt{python tasks.py run}}}}

\newcommand{\discussion}[3]{% {number}{question}{readings}
  \begin{frame}{Discussion #1}
    \begin{alertblock}{Question for the room}\large #2\end{alertblock}
    \vspace{2mm}{\small\textbf{Readings.} #3}
  \end{frame}}

\author{Aditya Sissodiya}
\institute{Lule\aa{} University of Technology \quad\textperiodcentered\quad Agentic AI Security, Identity, and Access Control (7.5 ECTS)}


\title{Can we prove it and measure it?}
\subtitle{Workshop 3 of 3 \textperiodcentered{} Safety analysis, evaluation, and what it all adds up to}
\date{\prep{set the workshop date}}

\begin{document}

\maketitle

\begin{frame}{Today}
  \begin{columns}[T]
    \column{0.57\textwidth}
    \small
    \begin{tabular}{@{}l l@{}}
      0:00 & Where Workshops 1 and 2 left us \\
      0:10 & Talk and live demos (40 min) \\
      0:50 & Hands-on: write a benchmark case \\
      1:20 & Break \\
      1:30 & Discussion 5: does decidability transfer? \\
      2:00 & Discussion 6: the smallest valid benchmark \\
      2:30 & Course synthesis and examiner questions \\
    \end{tabular}
    \column{0.39\textwidth}
    \begin{block}{Run along}
      \footnotesize\texttt{python tasks.py run w3}\\[2pt]
      Hands-on: \texttt{my\_cases.json} in \texttt{experiments/}
    \end{block}
  \end{columns}
\end{frame}

\begin{frame}{Learning outcomes for this workshop}
  \small After this workshop I should be able to show that I can:
  \begin{enumerate}
    \item[K3.1] explain the safety problem, why it is undecidable in general, and which restrictions make it decidable;
    \item[S3.2] model a delegation system as states and transitions and check safety properties by exhaustive search, with counterexamples;
    \item[S3.3] explain, with measured state counts, when a polynomial closure check suffices and when search is needed;
    \item[S3.4] design an access-control evaluation for agents with explicit ground truth and metrics, and run it on four designs;
    \item[J3.5] say what existing agent-security benchmarks measure and what they leave out;
    \item[J3.6] weigh security against utility and cost, state the limits of a simulation, and discuss the societal stakes of delegated authority.
  \end{enumerate}
\end{frame}

\begin{frame}{Where we are}
  \centering{\def\scenarioscale{0.75}% The running scenario used in all three workshops (agentsim/world.py).
% Optional scale: \def\scenarioscale{0.8} before \input.
\providecommand{\scenarioscale}{1}
\begin{tikzpicture}[scale=\scenarioscale, transform shape, >=Stealth,
  ent/.style={draw=black!60, rounded corners=2pt, fill=white, minimum height=8mm, minimum width=19mm,
              align=center, font=\scriptsize, inner sep=2pt},
  org/.style={draw=#1, fill=#1!7, rounded corners=4pt, inner sep=5pt, thick},
  orglabel/.style={font=\tiny\bfseries, text=#1},
  hop/.style={->, thick, black!75},
  side/.style={->, thick, dashed, #1}]
  \node[ent] (alice) {Alice\\[-1pt]\tiny operator};
  \node[ent, right=5mm of alice] (orch) {orchestrator\\[-1pt]\tiny agent};
  \node[ent, right=12mm of orch] (xrf) {XRF analyzer\\[-1pt]\tiny agent};
  \node[ent, right=12mm of xrf] (reg) {DPP registry\\[-1pt]\tiny MCP server};
  \node[ent, below=5mm of reg, minimum width=16mm] (p1) {battery-001\\[-1pt]\tiny owned here};
  \node[ent, below=2mm of p1, minimum width=16mm] (p2) {battery-002\\[-1pt]\tiny othermaker};
  \node[ent, below=11mm of xrf, draw=orgbroker] (harv) {harvester\\[-1pt]\tiny broker agent};
  \node[ent, below=11mm of orch] (inc) {incident agent\\[-1pt]\tiny regulator};
  \begin{scope}[on background layer]
    \node[org=orgrecycler, fit=(alice)(orch)] (recbox) {};
    \node[org=orglab, fit=(xrf)] (labbox) {};
    \node[org=orgmaker, fit=(reg)(p1)(p2)] (mkbox) {};
    \node[org=orgbroker, fit=(harv), dashed] (brbox) {};
    \node[org=orgregulator, fit=(inc)] (rgbox) {};
  \end{scope}
  \node[orglabel=orgrecycler, above=0pt of recbox] {recycler};
  \node[orglabel=orglab, above=0pt of labbox] {lab};
  \node[orglabel=orgmaker, above=0pt of mkbox] {manufacturer};
  \node[orglabel=orgbroker, below=0pt of brbox] {broker (adversary)};
  \node[orglabel=orgregulator, below=0pt of rgbox] {regulator};
  \draw[hop] (alice) -- node[above, font=\tiny]{1} (orch);
  \draw[hop] (orch) -- node[above, font=\tiny]{2} (xrf);
  \draw[hop] (xrf) -- node[above, font=\tiny]{3} (reg);
  \draw[black!50] (reg) -- (p1);
  \draw[side=orgbroker] (harv.east) -- (harv.east -| mkbox.west);
\end{tikzpicture}
}
  \begin{itemize}\small
    \item W1: the protocols leave cross-organization delegation to implementers; authority leaks at every hop.
    \item W2: ReBAC bounds the subject now, capabilities bound what was handed on; the hybrid meets six of seven requirements, provenance stays open.
    \item Today: can a design be shown safe, and how do we measure one?
  \end{itemize}
\end{frame}

\section{Proving it}

\begin{frame}{The safety question}
  \begin{columns}[T]
    \column{0.55\textwidth}
    \begin{itemize}\small
      \item Harrison, Ruzzo, and Ullman model protection as an access matrix changed by commands \ev{A}{harrison_protection_1976}.
      \item \emph{Safety}: can some right ever leak into a cell where it was not before?
      \item In general this is undecidable: commands can simulate a Turing machine.
      \item It becomes decidable under restrictions, for example when every command performs a single primitive operation.
    \end{itemize}
    \column{0.4\textwidth}
    \begin{itemize}\small
      \item The take-grant model decides whether a right can be shared in linear time \ev{A}{lipton_subject_1977}.
      \item Security analysis in trust management finds further decidable fragments \ev{A}{li_trustmanagement_2004}.
      \item A first formal framework for LLM-agent security: \ev{A}{formalizing_agentsecurity_2026}\,\pre.
    \end{itemize}
  \end{columns}
  \prep{Check each claim on this slide against notes R3.1--R3.3, including the exact restriction HRU prove decidable and its complexity. Resolve the Li and Tripunitara source (marked UNVERIFIED).}
\end{frame}

\begin{frame}{Our delegation system as a transition system}
  \begin{columns}[T]
    \column{0.5\textwidth}
    \small
    \textbf{State:} consent given?, revoked?, and a set of grants $(\mathit{holder}, \mathit{scope})$.\\[2pt]
    \textbf{Actions:} the owner consents (to Alice, scopes $S$, organizations $O$); a holder delegates; the owner revokes.\\[4pt]
    \textbf{Rule sets}
    \begin{itemize}\scriptsize
      \item[R0] anyone holding anything may delegate anything
      \item[R1] delegation must narrow
      \item[R2] + revocation cascades to derived grants
      \item[R3] + crossing into an organization needs the owner's consent
    \end{itemize}
    \column{0.46\textwidth}
    \small\textbf{Properties}
    \begin{itemize}\scriptsize
      \item[P1] no amplification: every grant's scope is in $S$
      \item[P2] after revocation, no grant survives
      \item[P3] every holder belongs to an organization in $O$
    \end{itemize}
    \vspace{2mm}
    \small Instance: Alice and her orchestrator (recycler), the XRF agent (lab), the broker's agent. $S = \{\text{materials}\}$, $O = \{\text{recycler}, \text{lab}\}$.
  \end{columns}
\end{frame}

\begin{frame}{Demo E3.1: four rule sets, three properties}
  \centering\result{safety_rules}\par
  \vspace{3mm}
  \result{safety_traces}
\end{frame}

\begin{frame}{Demo E3.2: what deciding it costs}
  \begin{columns}[c]
    \column{0.62\textwidth}
    \result{growth_plot}
    \column{0.36\textwidth}
    \small
    \begin{itemize}
      \item With a fixed set of principals the state space is finite, so search \emph{decides} each property, at exponential cost.
      \item Delegation only adds grants, so P1 and P3 can be decided on the closure: quadratic.
      \item P2 is about what happens after revocation; it needs search.
    \end{itemize}
  \end{columns}
\end{frame}

\begin{frame}{What this says about agents}
  \begin{itemize}
    \item A fixed cast of agents gives a finite, decidable problem. Agents that spawn sub-agents at run time bring back unbounded subject creation, the ingredient of HRU's undecidability.
    \item Properties that only forbid a bad grant can use cheap closure checks. Properties about revocation, expiry, or ordering need state, and get expensive.
    \item So a practical design restricts the rules (R3) until the properties it cares about are cheap to check, and checks the rest at run time.
  \end{itemize}
  \prep{Mark which of these three points are my inferences and which are established; test the first against R3.3 (formalizing agent security).}
\end{frame}

\section{Measuring it}

\begin{frame}{What existing benchmarks measure}
  \scriptsize
  \begin{tabularx}{\textwidth}{@{}l X c@{}}
    \toprule
    \textbf{Benchmark or defense} & \textbf{Measures} & \textbf{Multi-hop delegation ground truth?} \\
    \midrule
    AgentDojo \ev{A}{debenedetti_agentdojo_2024} & prompt-injection attacks and task utility for tool-using agents & \xmark \\
    $\tau$-bench \ev{A}{yao_taubench_2025} & policy adherence in tool-agent-user conversations & \xmark \\
    ToolEmu \ev{A}{ruan_toolemu_2024} & risky actions, with tools emulated by a language model & \xmark \\
    AgentHarm \ev{A}{andriushchenko_agentharm_2025} & compliance with harmful agentic requests & \xmark \\
    Agent Security Bench \ev{A}{zhang_asb_2025} & attacks and defenses across agent scenarios & \xmark \\
    OpenAgentSafety \ev{A}{vijayvargiya_openagentsafety_2026} & safety in realistic tool environments & \xmark \\
    Progent, MiniScope, CaMeL \ev{A}{progent_leastprivilege_2025}\ev{A}{miniscope_leastprivilege_2025}\ev{A}{camel_security_2025} & least-privilege or capability defenses, each with its own dataset & \pmark \\
    \bottomrule
  \end{tabularx}
  \prep{Recheck every row against the current releases and my notes R3.4--R3.12. The last column is my reading, not the authors' claim.}
\end{frame}

\begin{frame}{The evaluation design}
  \begin{columns}[T]
    \column{0.5\textwidth}
    \small\textbf{Ground truth:} allow only if
    \begin{enumerate}\scriptsize
      \item Alice is authorized for the field \emph{now};
      \item the field and passport are in the task;
      \item the caller is the delegated agent;
      \item nothing was revoked or expired;
    \end{enumerate}
    \small or if the emergency rule applies.\\[3pt]
    \small\textbf{14 cases in 9 categories}, written in JSON; anyone can add one.
    \column{0.46\textwidth}
    \small\textbf{Four designs}
    \begin{itemize}\scriptsize
      \item[B0] token passthrough
      \item[B1] OAuth: audience and scope, RFC 8693 delegation
      \item[B2] B1 + introspection, workload binding, ReBAC on the agent
      \item[B3] ReBAC on Alice at decision time + attenuated capability + cascading revocation
    \end{itemize}
    \small\textbf{Metrics:} unauthorized allows, false denies, injection contained, calls accepted after revocation, excess authority, simulated decision cost.
  \end{columns}
\end{frame}

\begin{frame}{Demo E3.3: results}
  \centering\result{bench_metrics}\par
  \vspace{2mm}
  {\scriptsize Excess authority: non-public (passport, field) pairs the delegated agent can read beyond the one it needs. Cost: simulated from assumed latencies (20 ms per introspection round trip), not measured.}
\end{frame}

\begin{frame}{Case by case}
  \begin{columns}[T]
    \column{0.55\textwidth}
    \result{bench_matrix}
    \column{0.42\textwidth}
    \small
    \begin{itemize}
      \item Every design rejects an untrusted issuer and an expired token. B1 adds only passthrough and requests beyond Alice's own rights.
      \item B2 stops theft and cross-tenant use, but not revocation at Alice's organization or Alice leaving.
      \item B3 gets every deny right and \alert{fails break-glass}: the strictest design has no emergency path.
    \end{itemize}
  \end{columns}
\end{frame}

\begin{frame}{Reading the results honestly}
  \begin{columns}[T]
    \column{0.48\textwidth}
    \textbf{Trade-offs}
    \begin{itemize}\small
      \item Security costs round trips: B3 contacts two authorization servers per decision.
      \item Least privilege fights emergency access. Break-glass needs its own semantics and an audit trail, not a looser policy.
      \item Revocation across organizations needs a channel the standards do not define.
    \end{itemize}
    \column{0.48\textwidth}
    \textbf{Limits of this evaluation}
    \begin{itemize}\small
      \item A simulation with a scripted agent; no LLM, no network.
      \item Costs are assumptions.
      \item I wrote the cases \emph{and} the designs: the ground truth may flatter B3.
      \item 14 cases show the shape of the problem; they are not a benchmark.
    \end{itemize}
  \end{columns}
\end{frame}

\section{Your turn}

\begin{frame}[fragile]{Hands-on: write a benchmark case (30 min)}
  \small
  \begin{enumerate}
    \item In pairs, find a situation where you expect B3 to be wrong.
    \item Write it as a case in \texttt{experiments/my\_cases.json} (see the activity sheet for the keys).
    \item Run all four designs on it:
\begin{verbatim}
python e3_benchmark.py --cases my_cases.json
\end{verbatim}
    \item If B3 gets it wrong: is the design wrong, or the ground truth?
  \end{enumerate}
\end{frame}

\discussion{5}{Do run-time composition and revocation push agent delegation outside the decidable fragments we know?}
{HRU \ev{A}{harrison_protection_1976}; take-grant \ev{A}{lipton_subject_1977}; formalizing LLM-agent security \ev{A}{formalizing_agentsecurity_2026}\,\pre. \prep{Prepare the strongest argument that it does not.}}

\discussion{6}{What is the smallest benchmark that would be valid for agent access control, and who should write its ground truth?}
{AgentDojo \ev{A}{debenedetti_agentdojo_2024}; Progent \ev{A}{progent_leastprivilege_2025}\,\pre; MiniScope \ev{A}{miniscope_leastprivilege_2025}\,\pre. \prep{Bring one concrete proposal for who owns the ground truth.}}

\section{What it adds up to}

\begin{frame}{Five things I would tell someone deploying agents across organizations}
  \begin{enumerate}
    \item Never pass tokens through. Exchange them, name the actor, and bind the token to the caller.
    \item Narrow authority to the task at every hop. Copying scope is how injected instructions become data leaks.
    \item Check the human principal at decision time, not only when the delegation was set up.
    \item Decide how revocation crosses organizations before the first incident, not after.
    \item Design the emergency path on purpose, with an audit trail, instead of leaving the strict policy to fail.
  \end{enumerate}
\end{frame}

\begin{frame}{Open problems}
  \small
  \begin{tabularx}{\textwidth}{@{}l X l@{}}
    \toprule
    \textbf{Problem} & \textbf{Why it is open} & \textbf{Seen in} \\
    \midrule
    Verifiable provenance & chains are either informational (RFC 8693) or self-asserted (caveats) & W2 R6 \\
    Cross-organization revocation & no standard channel; derived tokens outlive the origin & W1, W3 \\
    Break-glass semantics & strict designs deny emergencies; loose ones allow everything & W3 \\
    Decidable delegation with spawning agents & unbounded subject creation & W3 \\
    A shared benchmark & each defense brings its own dataset and ground truth & W3 \\
    Where governance lives & protocol, overlay, or registry & W1 \\
    \bottomrule
  \end{tabularx}
\end{frame}

\begin{frame}{Who carries the risk?}
  \begin{itemize}
    \item In public-safety and circular-economy settings, an agent acting on delegated authority can expose commercially sensitive data or delay an emergency response.
    \item When a delegated agent misuses authority, accountability depends on being able to say who authorized what. Today that record is optional.
    \item Strict designs move risk from data owners to responders; loose designs move it back.
  \end{itemize}
  \prep{Add one concrete case from the public-safety scenario and my position on who should decide the trade-off.}
\end{frame}

\begin{frame}{How this workshop was prepared}
  \centering\result{prep_summary}
\end{frame}

\begin{frame}[allowframebreaks]{References}
  \printbibliography[heading=none]
\end{frame}

\end{document}
